\documentclass[aspectratio=169]{beamer}

\usetheme{Madrid}
\usecolortheme{seahorse}
\setbeamertemplate{navigation symbols}{}

\usepackage[utf8]{inputenc}
\usepackage[T1]{fontenc}
\usepackage[spanish]{babel}
\usepackage{booktabs}
\usepackage{graphicx}
\usepackage{adjustbox}
\usepackage{amsmath}

\title{Development of a Heterogeneous Reconfigurable CGRA Cell}
\subtitle{Avance de resultados y analisis de desempeno}
\author{
Duarte Sanchez David \and
Jimenez Ulate Eddy \and
Ruiz Rivera Jeffrey \and
Esquivel Arias Fabian \and
Zuniga Ramirez Jose Eduardo
}
\institute{Escuela de Ingenieria en Electronica \\ Instituto Tecnologico de Costa Rica}
\date{2026}

\begin{document}

\begin{frame}
  \titlepage
\end{frame}

\begin{frame}{Motivacion}
\begin{itemize}
  \item Las CGRA ofrecen un punto intermedio entre flexibilidad (FPGA) y eficiencia (ASIC).
  \item En arreglos homogeneos hay subutilizacion de recursos cuando la carga no coincide con el tipo de PE.
  \item Propuesta: CGRA heterogenea con celdas especializadas para mejorar uso de recursos y desempeno.
\end{itemize}
\end{frame}

\begin{frame}{Objetivo del trabajo}
\begin{itemize}
  \item Disenar y validar una celda reconfigurable heterogenea para CGRA.
  \item Integrar:
  \begin{itemize}
    \item Celdas de ruteo
    \item Celdas de memoria distribuida
    \item Celda escalar
    \item Celda vectorial
    \item Celda MAC
  \end{itemize}
  \item Evaluar comportamiento funcional y metricas de rendimiento mediante SystemC + Vitis HLS.
\end{itemize}
\end{frame}

\begin{frame}{Arquitectura propuesta (2x2)}
\begin{itemize}
  \item Topologia base: malla $2\times2$ (escalable).
  \item Interconexion NoC con credit-based flow control.
  \item Interfaz tipo RISC-V (CSR/DMA) para configuracion y ejecucion.
  \item Enfoque modular para explorar variantes temporales y espaciales.
\end{itemize}
\end{frame}

\begin{frame}{Metodologia experimental}
\begin{itemize}
  \item Modelado y simulacion funcional en SystemC.
  \item Flujo HLS (Vitis 2024.1):
  \[
  csim \rightarrow csynth \rightarrow cosim \rightarrow export\_design
  \]
  \item Caso de estudio principal: \textbf{sum-reduction}.
  \item Comparacion entre mapeo:
  \begin{itemize}
    \item Temporal (1 celda MAC reutilizada)
    \item Espacial (4 celdas escalares)
  \end{itemize}
\end{itemize}
\end{frame}

\begin{frame}{Resultados funcionales preliminares}
\begin{itemize}
  \item Suite de regresion:
  \begin{itemize}
    \item 19 testbenches SystemC
    \item 139 aserciones
    \item 0 fallos funcionales
  \end{itemize}
  \item Sum-reduction end-to-end: \textbf{PASS}.
  \item Validacion cruzada HLS (csim): \textbf{PASS}.
\end{itemize}
\end{frame}

\begin{frame}{Uso de recursos (post-synthesis)}
\begin{table}[htbp]
\centering
\scriptsize
\begin{adjustbox}{max width=\linewidth}
\begin{tabular}{@{}lccc@{}}
\toprule
\textbf{Recurso} & \textbf{Temporal (MAC)} & \textbf{Espacial (Scalar)} & \textbf{Espacial/Temporal} \\
\midrule
LUT  & 1205 & 3857 & 3.20$\times$ \\
FF   & 1551 & 5230 & 3.37$\times$ \\
DSP  & 3    & 12   & 4.00$\times$ \\
SRL  & 0    & 20   & -- \\
BRAM & 0    & 0    & -- \\
\bottomrule
\end{tabular}
\end{adjustbox}
\end{table}
\end{frame}

\begin{frame}{Metricas especiales (sum-reduction)}
\begin{table}[htbp]
\centering
\scriptsize
\begin{adjustbox}{max width=\linewidth}
\begin{tabular}{@{}p{0.40\linewidth}ccc@{}}
\toprule
\textbf{Metrica} & \textbf{Temporal} & \textbf{Espacial} & \textbf{Espacial/Temporal} \\
\midrule
Operations per second (OP/s) & 19.63 MOP/s & 17.97 MOP/s & 0.92$\times$ \\
Latencia promedio por operacion & 50.95 ns/op & 55.65 ns/op & 1.09$\times$ \\
Tasa de utilizacion & 79.7\% & 47.8\% & 0.60$\times$ \\
\bottomrule
\end{tabular}
\end{adjustbox}
\end{table}
\end{frame}

\begin{frame}{Latencia medida en cosim}
\begin{itemize}
  \item Costo transaccion \texttt{start=true}:
  \begin{itemize}
    \item Temporal: 141 ciclos
    \item Espacial: 154 ciclos
  \end{itemize}
  \item Ambos cumplen timing con reloj logrado de 2.891 ns.
  \item Resultado clave:
  \begin{itemize}
    \item En este benchmark, temporal resulta ligeramente mas rapido y considerablemente menor en area.
  \end{itemize}
\end{itemize}
\end{frame}

\begin{frame}{Interpretacion tecnica}
\begin{itemize}
  \item El numero ideal de iteraciones de FSM no explica por si solo el rendimiento real.
  \item El costo por etapa/celda tras HLS afecta fuertemente la latencia final.
  \item Conclusiones para este caso:
  \begin{itemize}
    \item Temporal: mejor eficiencia de area y mejor tiempo total.
    \item Espacial: mayor paralelismo potencial, pero con mayor costo de recursos.
  \end{itemize}
\end{itemize}
\end{frame}

\begin{frame}{Estado de linea base RISC-V (gem5)}
\begin{itemize}
  \item Se prepara comparacion con baseline software RISC-V (RV64GC) en gem5 (SE mode).
  \item Modelos CPU considerados:
  \begin{itemize}
    \item Atomic
    \item Timing
    \item Minor
    \item O3
  \end{itemize}
  \item Metrica principal: ciclos (\texttt{numCycles}) e instrucciones (\texttt{numInsts}).
\end{itemize}
\end{frame}

\begin{frame}{Conclusiones}
\begin{itemize}
  \item La arquitectura heterogenea propuesta es funcional y reproducible.
  \item Sum-reduction confirma que especializacion de celda importa tanto como el estilo de mapeo.
  \item Para el caso evaluado, la version temporal con celda MAC ofrece mejor balance desempeno/area.
  \item Siguiente paso: ampliar benchmark y cerrar comparacion con baseline RISC-V en gem5.
\end{itemize}
\end{frame}

\begin{frame}{Trabajo futuro}
\begin{itemize}
  \item Escalar casos de prueba (tamano de problema y kernels).
  \item Profundizar en analisis de energia y area-delay product.
  \item Integrar resultados completos de gem5 para comparar HW/SW.
  \item Afinar mapeo y scheduling HLS por tipo de celda.
\end{itemize}
\end{frame}

\begin{frame}
\centering
\Huge Gracias
\end{frame}

\end{document}